% GENERATED FILE. Edit canonical lessons, then run npm run generate:books.
% canonical-source-hash: fnv1a64:d729c72a59fb1e4d
% canonical-lessons: GU-C304-tanavu, GU-C304-dubadvu, GU-C304-nuksan, GU-C304-aphat, GU-C304-sankat

\chapter{The Flood}
\label{ch:gu-the-flood}
\hlchaptermodality{304}

\begin{chapteropening}
\textbf{By the end of this chapter:} \emph{I can say what the flood carried away and what it cost.}

Complete the last lesson of chapter 304: I can say what the flood carried away and what it cost.
\end{chapteropening}

\section[\texorpdfstring{\gu{તણાવું} (taṇāvũ)}{તણાવું (taṇāvũ)}]{\gu{તણાવું} (taṇāvũ) — to be swept away}

\label{lesson:GU-C304-tanavu}



\begin{quote}
Before the new one: say the Gujarati for a hut, then the Gujarati for to lay waste.
\end{quote}

\subsection*{You'll want to know: \gu{તણાવું}}
\textbf{\gu{તણાવું}} — \emph{taṇāvũ} — \textquotedblleft{}to be swept away\textquotedblright{}.

\begin{quote}
\gu{પૂરમાં} \gu{અમારી} \gu{ગાય} \gu{તણાઈ} \gu{ગઈ}.
\end{quote}

\begin{quote}
\emph{pūrmā̃ amārī gāy taṇāī gaī.}
\end{quote}

\begin{quote}
Our cow was swept away in the flood.
\end{quote}

Past (it …): \gu{તે} \gu{ઝૂંપડું} \gu{તણાઈ} \gu{ગયું}, \emph{te jhū̃pḍũ taṇāī gayũ}.

\begin{cousinweb}
\textbf{\gu{તણાવું}} (\emph{taṇāvũ}) goes back to Sanskrit \emph{tan\'{\={o}}ti}, \textquotedblleft{}stretches\textquotedblright{}, and CDIAL lists the Gujarati verb, \textquotedblleft{}to be dragged\textquotedblright{} (CDIAL 5659): a flood drags things along with it. \textbf{\gu{તણાઈ} \gu{ગયું}}, it was swept away.
\end{cousinweb}

\begin{grammarlens}[title={What you've built}]
One more word for telling what happened.
\end{grammarlens}

\subsection*{Your turn}
\begin{itemize}
\raggedright
  \item \emph{Say it:} \emph{taṇāvũ}
  \item \emph{Say it:} \emph{taṇāvũ}, once more
  \item \emph{Say it:} \emph{pūrmā̃ amārī gāy taṇāī gaī.} — Our cow was swept away in the flood.
  \item \emph{Recall:} say the Gujarati for a hut, then the Gujarati for to lay waste, then say \emph{taṇāvũ} again
\end{itemize}

\begin{tcolorbox}[breakable,colback=teal!4,colframe=teal!35!black,title={Before you move on}]
What does \gu{તણાવું} mean? (\textquotedblleft{}To be swept away\textquotedblright{}.) Say it once more.
\end{tcolorbox}
\section[\texorpdfstring{\gu{ડુબાડવું} (ḍubāḍvũ)}{ડુબાડવું (ḍubāḍvũ)}]{\gu{ડુબાડવું} (ḍubāḍvũ) — to submerge; to flood}

\label{lesson:GU-C304-dubadvu}



\begin{quote}
Before the new one: say the Gujarati for to lay waste, then the Gujarati for to be swept away.
\end{quote}

\subsection*{You'll want to know: \gu{ડુબાડવું}}
\textbf{\gu{ડુબાડવું}} — \emph{ḍubāḍvũ} — \textquotedblleft{}to submerge; to flood\textquotedblright{}.

\begin{quote}
\gu{નદીએ} \gu{આખું} \gu{ગામ} \gu{ડુબાડી} \gu{દીધું}.
\end{quote}

\begin{quote}
\emph{nadīe ākhũ gām ḍubāḍī dīdhũ.}
\end{quote}

\begin{quote}
The river flooded the whole village.
\end{quote}

Past, with \gu{તેણે} (he or she …): \gu{તેણે} \gu{હોડી} \gu{ડુબાડી}, \emph{teṇe hoḍī ḍubāḍī}.

\begin{cousinweb}
\textbf{\gu{ડુબાડવું}} (\emph{ḍubāḍvũ}) is the causative of \textbf{\gu{ડૂબવું}}, to sink, from \emph{*ḍubb}, \textquotedblleft{}sink\textquotedblright{}, a form CDIAL reconstructs (CDIAL 5561): to make something go under. The long \emph{ū} of \textbf{\gu{ડૂબવું}} shortens in the longer verb.
\end{cousinweb}

\begin{grammarlens}[title={What you've built}]
One more word for telling what happened.
\end{grammarlens}

\subsection*{Your turn}
\begin{itemize}
\raggedright
  \item \emph{Say it:} \emph{ḍubāḍvũ}
  \item \emph{Say it:} \emph{ḍubāḍvũ}, once more
  \item \emph{Say it:} \emph{nadīe ākhũ gām ḍubāḍī dīdhũ.} — The river flooded the whole village.
  \item \emph{Recall:} say the Gujarati for to lay waste, then the Gujarati for to be swept away, then say \emph{ḍubāḍvũ} again
\end{itemize}

\begin{tcolorbox}[breakable,colback=teal!4,colframe=teal!35!black,title={Before you move on}]
What does \gu{ડુબાડવું} mean? (\textquotedblleft{}To submerge; to flood\textquotedblright{}.) Say it once more.
\end{tcolorbox}
\section[\texorpdfstring{\gu{નુકસાન} (nuksān)}{નુકસાન (nuksān)}]{\gu{નુકસાન} (nuksān) — damage}

\label{lesson:GU-C304-nuksan}



\begin{quote}
Before the new one: say the Gujarati for to be swept away, then the Gujarati for to submerge.
\end{quote}

\subsection*{You'll want to know: \gu{નુકસાન}}
\textbf{\gu{નુકસાન}} — \emph{nuksān} — \textquotedblleft{}damage\textquotedblright{}.

\begin{quote}
\gu{પૂરથી} \gu{ગામને} \gu{ખૂબ} \gu{નુકસાન} \gu{થયું}.
\end{quote}

\begin{quote}
\emph{pūrthī gāmne khūb nuksān thayũ.}
\end{quote}

\begin{quote}
The flood did the village a great deal of damage.
\end{quote}

A neuter noun.

\begin{cousinweb}
\textbf{\gu{નુકસાન}} (\emph{nuksān}) is usually traced, through Persian, to Arabic \emph{nuqṣān}, \textquotedblleft{}a decrease, a loss\textquotedblright{}. Harm is done with \textbf{\gu{થવું}}: \textbf{\gu{નુકસાન} \gu{થયું}}, there was damage.
\end{cousinweb}

\begin{grammarlens}[title={What you've built}]
One more word for telling what happened.
\end{grammarlens}

\subsection*{Your turn}
\begin{itemize}
\raggedright
  \item \emph{Say it:} \emph{nuksān}
  \item \emph{Say it:} \emph{nuksān}, once more
  \item \emph{Say it:} \emph{pūrthī gāmne khūb nuksān thayũ.} — The flood did the village a great deal of damage.
  \item \emph{Recall:} say the Gujarati for to be swept away, then the Gujarati for to submerge, then say \emph{nuksān} again
\end{itemize}

\begin{tcolorbox}[breakable,colback=teal!4,colframe=teal!35!black,title={Before you move on}]
What does \gu{નુકસાન} mean? (\textquotedblleft{}Damage\textquotedblright{}.) Say it once more.
\end{tcolorbox}
\section[\texorpdfstring{\gu{આફત} (āphat)}{આફત (āphat)}]{\gu{આફત} (āphat) — a calamity}

\label{lesson:GU-C304-aphat}



\begin{quote}
Before the new one: say the Gujarati for to submerge, then the Gujarati for damage.
\end{quote}

\subsection*{You'll want to know: \gu{આફત}}
\textbf{\gu{આફત}} — \emph{āphat} — \textquotedblleft{}a calamity\textquotedblright{}.

\begin{quote}
\gu{ગામના} \gu{લોકો} \gu{કહે} \gu{છે} \gu{કે} \gu{આ} \gu{સૌથી} \gu{મોટી} \gu{આફત} \gu{હતી}.
\end{quote}

\begin{quote}
\emph{gāmnā loko kahe che ke ā sauthī moṭī āphat hatī.}
\end{quote}

\begin{quote}
The villagers say that this was the worst calamity they have known.
\end{quote}

A feminine noun.

\begin{cousinweb}
\textbf{\gu{આફત}} (\emph{āphat}) is usually traced, through Persian, to Arabic \emph{āfa}, \textquotedblleft{}a blight, a calamity\textquotedblright{}. It is feminine: \textbf{\gu{મોટી} \gu{આફત}}, a great calamity; \textbf{\gu{આફત} \gu{આવી}}, disaster struck.
\end{cousinweb}

\begin{grammarlens}[title={What you've built}]
One more word for telling what happened.
\end{grammarlens}

\subsection*{Your turn}
\begin{itemize}
\raggedright
  \item \emph{Say it:} \emph{āphat}
  \item \emph{Say it:} \emph{āphat}, once more
  \item \emph{Say it:} \emph{gāmnā loko kahe che ke ā sauthī moṭī āphat hatī.} — The villagers say that this was the worst calamity they have known.
  \item \emph{Recall:} say the Gujarati for to submerge, then the Gujarati for damage, then say \emph{āphat} again
\end{itemize}

\begin{tcolorbox}[breakable,colback=teal!4,colframe=teal!35!black,title={Before you move on}]
What does \gu{આફત} mean? (\textquotedblleft{}A calamity\textquotedblright{}.) Say it once more.
\end{tcolorbox}
\section[\texorpdfstring{\gu{સંકટ} (sankaṭ)}{સંકટ (sankaṭ)}]{\gu{સંકટ} (sankaṭ) — a crisis; danger}

\label{lesson:GU-C304-sankat}



\begin{quote}
Before the new one: say the Gujarati for damage, then the Gujarati for a calamity.
\end{quote}

\subsection*{You'll want to know: \gu{સંકટ}}
\textbf{\gu{સંકટ}} — \emph{sankaṭ} — \textquotedblleft{}a crisis; danger\textquotedblright{}.

\begin{quote}
\gu{પૂર} \gu{પછી} \gu{ગામમાં} \gu{પાણીનું} \gu{સંકટ} \gu{હતું}.
\end{quote}

\begin{quote}
\emph{pūr pachī gāmmā̃ pāṇīnũ sankaṭ hatũ.}
\end{quote}

\begin{quote}
After the flood the village faced a water crisis.
\end{quote}

A neuter noun.

\begin{cousinweb}
\textbf{\gu{સંકટ}} (\emph{sankaṭ}) is a tatsama, Sanskrit \emph{saṁkaṭá}, \textquotedblleft{}narrow\textquotedblright{}, and as a noun \textquotedblleft{}a difficulty, a danger\textquotedblright{} (CDIAL 12817). It is neuter: \textbf{\gu{મોટું} \gu{સંકટ}}, a grave crisis.
\end{cousinweb}

\begin{grammarlens}[title={What you've built}]
One more word for telling what happened.
\end{grammarlens}

\subsection*{Your turn}
\begin{itemize}
\raggedright
  \item \emph{Say it:} \emph{sankaṭ}
  \item \emph{Say it:} \emph{sankaṭ}, once more
  \item \emph{Say it:} \emph{pūr pachī gāmmā̃ pāṇīnũ sankaṭ hatũ.} — After the flood the village faced a water crisis.
  \item \emph{Recall:} say the Gujarati for damage, then the Gujarati for a calamity, then say \emph{sankaṭ} again
\end{itemize}

\begin{tcolorbox}[breakable,colback=teal!4,colframe=teal!35!black,title={Before you move on}]
What does \gu{સંકટ} mean? (\textquotedblleft{}A crisis; danger\textquotedblright{}.) Say it once more.
\end{tcolorbox}
